\begin{table}[ht]
\caption{\textbf{\econ example}. The positive document describes the stochastic dominance, which is the underlying concept used to solve the problem.}
\centering
% \resizebox{\textwidth}{!}{
\begin{tabular}{p{13cm}}
\toprule
% \hspace{5.5cm} \econ \\
% \midrule
\textit{\textbf{Query}} \\
\midrule
Which of these two lotteries, a consumer with Von-Neumann Morgenstern preferences will choose under exponential distribution? \\
Consider two lotteries each having an exponential distribution. The function of cumulative distribution of an exponential distribution is: \\
$F(x;\lambda)=1-e^{-\lambda x} \forall x \in R_+.$ \\
The expected gain given by this distribution is $E[X]=1/\lambda$. \\
Suppose that the first lottery, $F_1$ has a parameter $\lambda_1$ and the second, $F_2$ has a parameter $\lambda_2$. Suppose that $\lambda_1$<$\lambda_2$. Which of these two lotteries, a consumer with Von-Neumann Morgenstern preferences will choose? \\
I am confused on this question. \\
\midrule
\textit{\textbf{Chain-of-thought reasoning to find documents}} \\
\midrule
The expected value is different from the expected utility. The preferences under uncertainty in this case should be linked to first-order stochastic dominance. We need to look up the detailed definition to understand how the problem could be solved by using that framework. \\
\midrule
\textit{\textbf{Example positive document}} \\
\midrule
Statewise dominance implies first-order stochastic dominance (FSD), which is defined as: \\
Random variable A has first-order stochastic dominance over random variable B if for any outcome x, A gives at least as high a probability of receiving at least x as does B, and for some x, A gives a higher probability of receiving at least x. In notation form, $P[A\geq x] \geq P[B\geq x]$ for all x, and for some x, $P[A\geq x] > P[B\geq x]$. \\
In terms of the cumulative distribution functions of the two random variables, A dominating B means that $F_A(x) \leq F_B(x)$ for all x, with strict inequality at some x. \\
In the case of non-intersecting distribution functions, the Wilcoxon rank-sum test tests for first-order stochastic dominance. \\
... \\
Extended example \\
Consider three gambles over a single toss of a fair six-sided die:

State (die result)
1
2
3
4
5
6 \\
gamble A wins 
\$
1
1
2
2
2
2 \\
gamble B wins 
\$
1
1
1
2
2
2 \\
gamble C wins 
\$
3
3
3
1
1
1 \\
Gamble A statewise dominates gamble B because A gives at least as good a yield in all possible states (outcomes of the die roll) and gives a strictly better yield in one of them (state 3). Since A statewise dominates B, it also first-order dominates B. \\
... \\
\midrule
\textit{\textbf{Example negative document}} \\
\midrule
Which Lottery Is Easiest To Win Within the U.S.? \\
Winning the lottery is one of the least likely things to happen in your life. However, some lotteries are much easier to win than others. \\
It’s common for a lot of Americans to daydream about winning the lottery. These dreams became a reality for a few lucky winners as they became millionaires overnight. The cold truth is that the overwhelming majority of lottery dreams will remain dreams.

Most people know that the odds of winning the lottery are not in your favor. If it were easy to win, everyone would play, and the prizes would only be a few cents. 

You’re likelier to be struck by lightning, live to be 110 years old, and be declared a saint by the Catholic Church. 

The combination of high jackpots and low odds has caused many people to try and cheat in the lottery. The attempts might be successful for a short period, but they always end with prison sentences. 

Fortunately, there are a few ways that you can legally improve your odds of winning. It will start by choosing games that have the best odds. 

... \\
\bottomrule
\end{tabular}
% }
\label{tab:economic_example}
\end{table}